% =====================================================================
% Figure contract
% 核心结论：CUA 只有把执行、核验恢复与外层学习闭合为可追责状态转移，才不会把未生效动作误判为成功。
% 图原型：参考 cua-survey-fig1-overview.tex 中央 hero 的扩展版；中等复杂度纯结构主图。
% 证据层级：hero = 三色执行闭环；支撑 = 路径图例与两条失败统计注释。
% 能不能更少？：不能；11 个节点、14 条边、三色图例和两条注释均由本图规格逐项要求。
%
% Step ① 设计文档（Form A）
% 领域/格式：Computer-Use Agents 中文综述；TikZ standalone + XeLaTeX。
% 印刷：W_print = 22 cm（整页横排）；目标最小字号 = 6.5 pt；设计画布 W_ink <= 24 cm。
% 档位/流向：中等复杂度；蓝色主流左→右，暖色失败流向左回折，绿色学习流由底层写回。
% 模块：主行 6 个；V/R 第二层；H 第三层；D/L 底层；左下图例；底部两行证据注释。
% 空间：H→S 走最外右/上 rail；R→P 走主行下方暖色 rail；L→O/P 共享绿色 trunk 后分成两个箭头。
% 视觉强调：所有语义节点统一中性描边，用 pastel fill 区分角色；路径色只编码三类边。
% ASCII 草图：
%        title          H→S warm dashed outer/top rail
%   [I]→[O]→[S]→[P]→[A]→[E]
%    legend          ↑warm   ↓warm
%                    └──[R]←[V]
%                         ↓dash        green V→L bottom rail
%                        [H]      [D]→[L]
%                                  └green→O / →P
%   note: false completion
%   note: repeated invalid actions
% 预防：三条暖色回环分轨；两条学习写回分轨；所有跨层线用 named waypoint 正交绕障。
%
% Edge contract (exactly 14 semantic edges):
%  1 I→O  2 O→S  3 S→P  4 P→A  5 A→E
%  6 E→V  7 V→R  8 R→P  9 R→H  10 H→S
% 11 V→L 12 D→L 13 L→O 14 L→P
% =====================================================================

\documentclass[border=8pt]{standalone}
\usepackage{fontspec}
\usepackage{xeCJK}
\usepackage{xcolor}
\usepackage{tikz}

\setCJKmainfont{Heiti SC}
\setCJKsansfont{Heiti SC}
\setsansfont{Helvetica Neue}

\usetikzlibrary{arrows.meta,bending,positioning,calc,fit,backgrounds}

% Low-saturation academic palette. Edge colors are the semantic encoding.
\definecolor{ink}{HTML}{293747}
\definecolor{muted}{HTML}{737E8A}
\definecolor{cardBorder}{HTML}{6C7986}
\definecolor{forwardEdge}{HTML}{5F7C96}
\definecolor{recoveryEdge}{HTML}{9A7064}
\definecolor{learningEdge}{HTML}{69866D}

\definecolor{inputFill}{HTML}{EAF1F6}
\definecolor{stateFill}{HTML}{EFEAF3}
\definecolor{policyFill}{HTML}{F4F0E5}
\definecolor{actionFill}{HTML}{F4E9EA}
\definecolor{envFill}{HTML}{EDF0F2}
\definecolor{verifyFill}{HTML}{E6F0EF}
\definecolor{recoveryFill}{HTML}{F4EAE7}
\definecolor{humanFill}{HTML}{F3E9EC}
\definecolor{learnFill}{HTML}{EAF1E8}
\definecolor{factoryFill}{HTML}{F3EFE3}
\definecolor{legendFill}{HTML}{F7F8F9}

\pgfdeclarelayer{background}
\pgfsetlayers{background,main}

\tikzset{
  every node/.style={
    font=\sffamily\fontsize{8.7}{10.4}\selectfont,
    text=ink
  },
  card/.style={
    rectangle,
    rounded corners=5pt,
    draw=cardBorder,
    line width=0.75pt,
    minimum width=3.10cm,
    minimum height=1.35cm,
    text width=2.72cm,
    align=center,
    inner xsep=5pt,
    inner ysep=5pt
  },
  title/.style={
    font=\sffamily\fontsize{13.0}{15.5}\selectfont\bfseries,
    text=ink,
    align=center
  },
  edge label/.style={
    font=\sffamily\fontsize{7.6}{9.0}\selectfont,
    text=recoveryEdge,
    fill=white,
    fill opacity=0.92,
    text opacity=1,
    inner xsep=3pt,
    inner ysep=1.5pt,
    rounded corners=2pt
  },
  legend box/.style={
    rectangle,
    rounded corners=4pt,
    draw=cardBorder!55,
    fill=legendFill,
    line width=0.55pt,
    minimum width=6.15cm,
    minimum height=2.70cm,
    inner sep=0pt
  },
  legend title/.style={
    font=\sffamily\fontsize{8.5}{10.2}\selectfont\bfseries,
    text=ink,
    anchor=west
  },
  legend text/.style={
    font=\sffamily\fontsize{7.8}{9.2}\selectfont,
    text=muted,
    anchor=west
  },
  note/.style={
    font=\sffamily\fontsize{7.25}{8.9}\selectfont,
    text=muted,
    align=left,
    inner sep=0pt
  },
  % Canonical long and short arrow tips. No edge style adds rounded corners.
  forward/.style={
    -{Stealth[length=5pt 1.5,width'=0pt 0.6,sep=0pt 0.5]},
    line width=1.15pt,
    shorten >=2pt,
    shorten <=1pt,
    color=forwardEdge
  },
  forward short/.style={
    -{Stealth[length=3pt 1.0,width'=0pt 0.5,sep=0pt 0.3]},
    line width=0.9pt,
    shorten >=1pt,
    shorten <=0pt,
    color=forwardEdge
  },
  warm/.style={
    -{Stealth[length=5pt 1.5,width'=0pt 0.6,sep=0pt 0.5]},
    line width=1.15pt,
    shorten >=2pt,
    shorten <=1pt,
    color=recoveryEdge
  },
  warm short/.style={
    -{Stealth[length=3pt 1.0,width'=0pt 0.5,sep=0pt 0.3]},
    line width=0.9pt,
    shorten >=1pt,
    shorten <=0pt,
    color=recoveryEdge
  },
  warm dashed short/.style={
    warm short,
    densely dashed
  },
  green/.style={
    -{Stealth[length=5pt 1.5,width'=0pt 0.6,sep=0pt 0.5]},
    line width=1.15pt,
    shorten >=2pt,
    shorten <=1pt,
    color=learningEdge
  },
  green short/.style={
    -{Stealth[length=3pt 1.0,width'=0pt 0.5,sep=0pt 0.3]},
    line width=0.9pt,
    shorten >=1pt,
    shorten <=0pt,
    color=learningEdge
  },
  legend forward/.style={line width=1.15pt,color=forwardEdge},
  legend warm/.style={line width=1.15pt,color=recoveryEdge},
  legend green/.style={line width=1.15pt,color=learningEdge}
}

\begin{document}
\sffamily
\begin{tikzpicture}

% ---------------------------------------------------------------------
% Eleven semantic cards, positioned by construction.
% ---------------------------------------------------------------------
\node[card,fill=inputFill] (intent) {用户意图\\[-1pt]{\bfseries I}};
\node[card,fill=inputFill,right=0.55cm of intent] (observation)
  {Observation·Belief\\[-1pt]Source {\bfseries O}};
\node[card,fill=stateFill,right=0.55cm of observation] (state)
  {Explicit Task State\\[-1pt]{\bfseries S}};
\node[card,fill=policyFill,right=0.55cm of state] (plan)
  {Planning·Policy\\[-1pt]{\bfseries P}};
\node[card,fill=actionFill,right=0.55cm of plan] (action)
  {Semantic GUI·API\\[-1pt]Action {\bfseries A}};
\node[card,fill=envFill,right=0.55cm of action] (environment)
  {Environment\\[-1pt]Transition {\bfseries E}};

\node[card,fill=verifyFill,below=1.35cm of environment] (verifier)
  {Verifier·Feedback\\[-1pt]{\bfseries V}};
\node[card,fill=recoveryFill,left=0.55cm of verifier] (recovery)
  {Recovery·Abstention\\[-1pt]{\bfseries R}};
\node[card,fill=humanFill,below=0.90cm of recovery] (human)
  {Human Handoff\\[-1pt]{\bfseries H}};

\node[card,fill=learnFill,below=3.95cm of plan] (learning)
  {Learning\\[-1pt]{\bfseries L}};
\node[card,fill=factoryFill,left=0.75cm of learning] (factory)
  {Data·Task Factory\\[-1pt]{\bfseries D}};

% Title: centered over the six-card forward path.
\coordinate (mainTopCenter) at ($(intent.north west)!0.5!(environment.north east)$);
\node[title,anchor=south] at ([yshift=1.12cm]mainTopCenter |- intent.north)
  {CUA 执行闭环（§3.2）：前向路径、核验返回与外层学习};

% ---------------------------------------------------------------------
% Exact 14-edge contract.
% ---------------------------------------------------------------------

% 1--5. Forward execution: solid blue, all adjacent gaps < 1.5 cm.
\draw[forward short] (intent.east) -- (observation.west);       % 1 I→O
\draw[forward short] (observation.east) -- (state.west);       % 2 O→S
\draw[forward short] (state.east) -- (plan.west);              % 3 S→P
\draw[forward short] (plan.east) -- (action.west);             % 4 P→A
\draw[forward short] (action.east) -- (environment.west);      % 5 A→E

% 6--8. Verification and recovery: solid warm.
\draw[warm short] (environment.south) -- (verifier.north);     % 6 E→V
\draw[warm short] (verifier.west) -- (recovery.east);          % 7 V→R

% R→P uses a dedicated rail below the main row; it never overlaps a node.
\coordinate (recoveryRailStart) at ([yshift=0.55cm]recovery.north);
\coordinate (planWarmTarget) at ([xshift=0.88cm]plan.south);
\coordinate (recoveryRailTurn) at (planWarmTarget |- recoveryRailStart);
\draw[warm short]
  (recovery.north) -- (recoveryRailStart)
  -- (recoveryRailTurn) -- (planWarmTarget);                    % 8 R→P

% 9--10. Human intervention: dashed warm.
\draw[warm dashed short] (recovery.south) -- (human.north);    % 9 R→H

% H→S takes the outer-right and top rail. The label sits on its longest span.
\coordinate (handoffRailRight) at
  ([xshift=0.65cm,yshift=0.60cm]environment.north east);
\coordinate (handoffFromHuman) at (human.east -| handoffRailRight);
\coordinate (handoffStateTop) at (state.north |- handoffRailRight);
\draw[warm dashed short]
  (human.east) -- (handoffFromHuman) -- (handoffRailRight)
  -- node[edge label,below=3pt] {介入信息写回 state}
  (handoffStateTop) -- (state.north);                            % 10 H→S

% 11. V→L uses a bottom rail, below H and both learning cards.
\coordinate (learnBottomRailV) at ([yshift=-0.62cm]learning.south -| verifier.south);
\coordinate (learnBottomRailL) at (learning.south |- learnBottomRailV);
\draw[green short]
  (verifier.south) -- (learnBottomRailV)
  -- (learnBottomRailL) -- (learning.south);                     % 11 V→L

% 12. Data/task factory directly feeds learning.
\draw[green short] (factory.east) -- (learning.west);           % 12 D→L

% 13--14. L→O/P: a green trunk rises through the S/P gap, then splits.
% The spine sits above the legend and to the left of the warm R→P endpoint.
\coordinate (learnObsTarget) at ([xshift=0.18cm]observation.south);
\coordinate (learnPlanTarget) at (plan.south);
\coordinate (learnJunction) at ([yshift=-0.98cm]plan.south);
\coordinate (learnObsRailTurn) at (learnObsTarget |- learnJunction);
\draw[line width=0.9pt,color=learningEdge]
  (learning.north) -- (learnJunction);
\draw[green short]
  (learnJunction) -- (learnObsRailTurn) -- (learnObsTarget);     % 13 L→O
\draw[green short]
  (learnJunction) -- (learnPlanTarget);                          % 14 L→P

% ---------------------------------------------------------------------
% Three-color path legend: swatches reuse the exact edge color names.
% ---------------------------------------------------------------------
\node[legend box,anchor=north west] (legend)
  at ([yshift=-1.20cm]intent.south west) {};
\node[legend title] at ([xshift=0.35cm,yshift=-0.38cm]legend.north west)
  {路径图例};

\draw[legend forward]
  ([xshift=0.38cm,yshift=-1.00cm]legend.north west) -- ++(1.10cm,0);
\node[legend text] at ([xshift=1.70cm,yshift=-1.00cm]legend.north west)
  {前向执行};

\draw[legend warm]
  ([xshift=0.38cm,yshift=-1.60cm]legend.north west) -- ++(1.10cm,0);
\node[legend text] at ([xshift=1.70cm,yshift=-1.60cm]legend.north west)
  {核验-恢复-介入};

\draw[legend green]
  ([xshift=0.38cm,yshift=-2.20cm]legend.north west) -- ++(1.10cm,0);
\node[legend text] at ([xshift=1.70cm,yshift=-2.20cm]legend.north west)
  {外层学习};

% ---------------------------------------------------------------------
% Evidence notes. The explicit width keeps both lines inside the 24 cm canvas.
% ---------------------------------------------------------------------
\coordinate (notesAnchor) at (intent.west |- factory.south);
\node[note,anchor=north west,text width=21.35cm] (noteOne)
  at ([yshift=-1.48cm]notesAnchor)
  {失败任务中 \textgreater86\% 是 false completion——动作没生效却继续相信成功，即 V 未核验、R 未触发（VLAA-GUI）};
\node[note,anchor=north west,text width=21.35cm]
  (noteTwo) at ([yshift=-0.43cm]noteOne.south west)
  {72.3\% 的失败来自重复无效动作导致 timeout（VeriGUI）};

\end{tikzpicture}
\end{document}
